\chapter{Cognitive Deficits (Brain Fog and Memory Lapses)}

\section*{Definition}
Cognitive deficits refer to impairments in mental processes such as memory, attention, executive function, language, and visuospatial skills. These deficits can range from mild subjective complaints to severe functional impairment as in dementia.

\section*{When to See a Doctor}
\begin{itemize}
\item Memory loss that disrupts daily life
\item Difficulty with planning, problem-solving, or multitasking
\item Word-finding difficulty or trouble following conversations
\item Getting lost in familiar places
\item Changes in judgment or decision-making
\item Personality changes, apathy, or social withdrawal
\item Difficulty with familiar tasks (cooking, finances)
\item Onset of confusion or delirium
\end{itemize}

\section*{How It Develops}
Cognitive function depends on complex neural networks distributed across the brain. Deficits arise when these networks are disrupted by structural damage (stroke, trauma), neurodegenerative processes (Alzheimer disease), metabolic derangements, neuroinflammation, or psychiatric conditions. Cholinergic circuits in the basal forebrain and hippocampus are particularly vulnerable.

\section*{Causes}
\begin{itemize}
\item Neurodegenerative diseases (Alzheimer disease, frontotemporal dementia)
\item Cerebrovascular disease (stroke, vascular dementia)
\item Traumatic brain injury
\item Mild cognitive impairment (MCI)
\item Depression, anxiety, or chronic stress (pseudodementia)
\item Sleep disorders (sleep apnea, insomnia)
\item Vitamin B12 deficiency
\item Thyroid dysfunction
\item Medication side effects (anticholinergics, benzodiazepines, opioids)
\item Alcohol or substance abuse
\item Chronic medical conditions (renal failure, hepatic encephalopathy)
\end{itemize}

\section*{Related Symptoms}
\begin{itemize}
\item Memory loss
\item Confusion
\item Disorientation
\item Difficulty concentrating
\item Fatigue
\item Mood changes (depression, anxiety, irritability)
\item Sleep disturbances
\item Headache
\end{itemize}

\section*{Medical Evaluation}
\begin{itemize}
\item Detailed history from patient and family member
\item Cognitive screening (Montreal Cognitive Assessment)
\item Neuropsychological testing
 \item Brain imaging (usually MRI or CT) when indicated to assess structural or vascular causes
 \item Targeted blood tests such as vitamin B12, thyroid testing, and a complete blood count when clinically indicated
\item Assessment for depression
 \item Lumbar puncture or other specialized testing only in selected cases, such as suspected inflammatory, infectious, or rapidly progressive disease
\end{itemize}

\section*{Treatment Options}
\begin{itemize}
\item Treat reversible causes: replace B12, correct thyroid function
 \item Cholinesterase inhibitors or memantine may be considered for selected stages and types of dementia; they are not treatments for every cause of cognitive symptoms
\item Antidepressants for depression-related cognitive impairment
\item Cognitive rehabilitation and occupational therapy
\item Management of vascular risk factors
\item Referral to a memory clinic or neurologist
\end{itemize}
\section*{Self-Care and Home Management}
\begin{itemize}
\item Engage in regular mental stimulation (reading, puzzles)
\item Maintain a consistent daily routine and use memory aids
\item Get adequate sleep (7-9 hours per night)
\item Exercise regularly to improve cerebral blood flow
 \item Follow a balanced dietary pattern such as a Mediterranean-style diet; no supplement has been proven to prevent or reverse cognitive impairment in everyone
\item Reduce alcohol consumption
\item Stay socially connected
\item Manage chronic conditions (hypertension, diabetes, depression)
\end{itemize}

\section*{References}
\begin{thebibliography}{99}
\bibitem{cognitive_deficits:ref1} Livingston G, Huntley J, Sommerlad A, et al. ``Dementia Prevention, Intervention, and Care: 2020 Report of the Lancet Commission.'' Lancet. 2020;396(10248):413--446.
\bibitem{cognitive_deficits:ref2} McKhann GM, Knopman DS, Chertkow H, et al. ``The Diagnosis of Dementia Due to Alzheimer's Disease: Recommendations from the NIA-AA Workgroups.'' Alzheimers Dement. 2011;7(3):263--269.
\bibitem{cognitive_deficits:ref3} Petersen RC. ``Mild Cognitive Impairment.'' N Engl J Med. 2011;364(23):2227--2234.
\bibitem{cognitive_deficits:ref4} Hugo J, Ganguli M. ``Dementia and Cognitive Impairment: Epidemiology, Diagnosis, and Treatment.'' Clin Geriatr Med. 2014;30(3):421--442.
\bibitem{cognitive_deficits:ref5} Feldman HH, Jacova C, Robillard A, et al. ``Diagnosis and Treatment of Dementia: 2. Diagnosis.'' CMAJ. 2008;178(7):825--836.
\bibitem{cognitive_deficits:ref6} Knopman DS, Boeve BF, Petersen RC. ``Essential Neurology for the Non-Neurologist.'' 2nd ed. Lippincott Williams \& Wilkins; 2005.
\end{thebibliography}
